\section{A multiplicity-scaled dyadic lower profile at early normalized source ranks}
Put $c_0=1$, $c_n=2^{n-1}$ for $n\ge1$. Thus
$\sum_{i=0}^n c_i=2^n$ and $\sum_{i<n}c_i=c_n$ for $n\ge1$.

\begin{theorem}[Early normalized source lower bounds]
For every actual increasing arm, write $\kappa\ge2$ for its first
missing weight, $\delta=E(\kappa)$, and $A=z^{-\kappa}E$.
Then
\[
 A_n\ge c_n\quad(0\le n\le\kappa-2),                 \tag{1}
\]
\[
 A_n\ge\delta c_n\quad(0\le n\le\kappa-\delta).       \tag{2}
\]
The same conclusions hold throughout the formal append grammar
with nondecreasing first missing weights. These are coefficient
lower bounds in the stated rank windows, not global dominance of
an infinite dyadic series or dominance after differentiation.
\end{theorem}
\begin{proof}
For an initial all-ones arm, $\kappa=2$. If $\delta=2$, $A_0=2$
and (2) only asks about zero. If $\delta=1$, $A_0=1$ and $A_1\ge1$,
which are exactly its required two ranks. This also proves (1)'s base.
Induct using the three exact normalized source recurrences.

First prove (1). In a growth step from old first weight $k$ to
$k+1$, $A'=[t]A+z^{t-k}$ with $t\ge k$. At a new required rank
$1\le n\le k-1$, the window contains ALL old ranks $0$ through
$n-1$, each within the old bound $k-2$. Their sum is at least $c_n$.
At zero the old positive constant supplies one. In a stationary
step $t=k-1$, $A'=1+z[t]A$; at every $1\le n\le k-2$ its
coefficient is the same full prefix through $n-1$, proving (1).

For (2), a growth step preserving $\delta=d$ enlarges its required
rank range by one. At any new $n\ge1$ in that range, the old
scaled bound applies through $n-1$, and the complete window prefix
has sum at least $dc_n$. The unchanged constant handles zero.
In a tied growth, $k$ and $d$ both increase by one, so the required
range $n\le k-d$ is unchanged. At a nonzero such rank, all old
coefficients through $n$ are covered and give at least $d2^n$.
Since $d\ge1$, this is at least $(d+1)2^{n-1}=(d+1)c_n$.
The new constant is exactly $d+1$.
Finally a stationary step resets the multiplicity to one. For its
entire new range $1\le n\le k-1$, the unscaled bound (1) covers
all old ranks through $n-1$, whose sum is $c_n$.
This proves the simultaneous unbounded induction.
\end{proof}

\begin{corollary}[A first-overlap source payment when the larger weight is far away]
At the first overlap orient $3\le p\le q$, with source
multiplicities $d\le p$, $e\le q$. Assume both forward primitives
strictly increase there. If $q\ge p+e+1$, then
\[
               \Delta E_Z(p+q)+\Delta E_Y(p+q)\ge de.
\]
Thus this first-overlap S branch is paid, without C. The usual
OLD A2 head input and nonnegative tail close all heads; zero head
requires neither C nor OLD A2 as in the previous overlap proofs.
\end{corollary}
\begin{proof}
If $e\ge2$, the last step cannot be stationary. In a tied growth
its prefix has first weight $q-1$ and multiplicity $e-1$, so (2)
covers relative rank $p$. The uncut exact difference is at least
$(e-1)2^{p-1}$. In a preserving growth, (2) covers the same rank
and gives $e2^{p-1}$, with a possible unit dent. This again is at
least $(e-1)2^{p-1}$. For $p\ge4$, $2^{p-1}\ge2p$ and
$2(e-1)\ge e$, so this pays $de\le pe$.
If $p=3$, strict increase of $D_Z$ through rank three gives
$2-d=D_Z(3)\ge D_Z(2)=1$, hence $d=1$.
Now $4(e-1)\ge e$ pays $de=e$.
If $e=1$, the last step may be stationary. A stationary step gives
source increment equal to a prefix's normalized coefficient at
rank $p-1$, which by (1) is at least $2^{p-2}\ge1$.
For a preserving growth, (2) gives source increment at least
$2^{p-1}-1\ge p\ge d$.
In the stationary case, the delayed source on the other factor
also contributes $D_Z(q)\ge D_Z(p)=2^{p-2}-d$.
The sum is at least $2^{p-1}-d\ge d$ for $p\ge4$ and
$d\le\min(p,2^{p-3})$; at $p=3$, $d=1$ and both delayed
source increments are at least one. Thus the total source payment
covers $de$ in every case.
\end{proof}

Combining this far-weight branch with the completed large-multiplicity
profile payment reduces the still-unproved first-overlap cases to
finitely many first weights and multiplicities. This reduction does
not itself prove their required source minima or any later overlap.
